\clearpage
\section*{2026-09-15: A twelve-parent review finds substantial land-area discrepancies}
\textit{Agent-drafted research entry.}

Ten provisional footprint comparisons give a mean land area of 32,755.5 square
feet, compared with 50,076.6 in production: a reduction of 17,321.1 square feet
per parent, or 34.6\% of the original mean. Seven areas decrease and three
increase. The mean absolute difference is 23,513.7 square feet and the median
absolute percentage difference is 55.4\%. These magnitudes show that the
flagged records can contain substantial measurement problems.

The sample was drawn before review using seed 20260915 after sorting the 166
flagged, composition-eligible parents by ID: nine historical and three
post-policy parents. Eleven of the twelve have partial lot-change flags.
Two remain unresolved, leaving eight historical and two post-policy comparisons.
Each compared parent counts once. The small sample, stratified draw, and
exclusion of unresolved cases preclude treating this as an estimate for all
166 parents. The median signed percentage change is $-26.4\%$. Averaging the
individual percentage changes instead gives $+53.6\%$, driven by two projects
with very small initial recorded areas; this differs from the change in the mean.

The comparison sums published PLUTO 25v4 lot areas for the mapped filing BBLs.
Dated DOF maps and transaction records explain how those parcels relate to the
earlier lots. Across the ten usable comparisons, PLUTO and the September 15
DOF geometry agree on more than 99.99\% of their union. The resulting footprints
are plausible candidates, not adopted measurements of land available at the
initial filing date. Production characteristics, weights, and model fits remain
provisional and unchanged by this review.

At 11 Gerry Street, production assigns all 111,040 square feet of a predecessor
lot; the mapped filing parcel has 18,044 square feet. At 1166 Fox Street,
production retains an earlier 2,500-square-foot parcel, while DOF documents a
merger and reapportionment leaving a 10,015-square-foot filing parcel. Weirfield
similarly reuses a lot number after changing its boundaries: 3,050 square feet
in production versus 19,341 in PLUTO 25v4. A matching identifier can therefore
refer to different land at different dates.

Sackett and Noble/Oak are omitted from the averages. Sackett's filing lot 47
has 12,500 square feet on the September 2025 map, but the March 2026 NYSDEC
amendment describes a site containing lots 17 and 47 together. Lot 47 merges
into 17 in August 2026, after the July 8 endpoint. Its original project extent
needs reconciliation. Noble's proposed lot 10 is unmapped, and lot 1 has a
waterfront boundary discrepancy. Godwin also illustrates an attribute problem:
allocating the earlier lot's recorded area by geometric overlap gives about
8,075 square feet, whereas PLUTO records 12,541 for the mapped new parcel.
The diagnostic allocation is retained separately from the comparison area.

\small
\textit{Reproduction.} Run \texttt{make dof-subset-review}. The existing
\path{tasks/audits/audit_hdb_mappluto_condo_recovery/} task produces
a parent comparison table, an average-change table, and a twelve-page footprint
PDF. Its small review table records the dated
map IDs, transaction IDs, and documentary judgments. The acquisition task
preserves the original geometry response and checksummed download rules for
25 dated maps and the NYSDEC amendment. The detailed findings and formulas are
in \texttt{report/dof\_subset\_review.md}.
\normalsize


\textit{September 21 update.} The current rebuild includes Godwin's accepted
companion lot 88 and uses its merged 18,441-square-foot production footprint.
A separately dated DOF geometry response supplies lot 88; the original geometry
snapshot and assessment bytes remain unchanged. The numerical comparison above
records the original September 15 draw. Current output tables compare the
same selected IDs with the updated parent definition and production areas.
